\chapter{量子ゆらぎと構造の起源}

\lectureoverview{最終講義では、インフレーションによって滑らかにされた宇宙が、それでもなお銀河が成長するもととなる微小なむらをどのように獲得できるのかを問う。答えは段階的に組み立てられる。まず普遍定数がプランクスケールを定義し、身近な素粒子物理学と量子重力との間の途方もない隔たりを明らかにする。次にマイクロ波の空が目標を与える。すなわち、\(10^{-5}\)程度の温度と密度の変動が多数のスケールに広がっている。零点運動は避けられないゆらぎを与え、減衰は振動がいかに凍結しうるかを説明し、膨張宇宙の波動方程式はインフラトンの各フーリエモードをまさにそのような減衰振動子に変える。モードはその物理的波長がハッブルスケールに達したときに凍結する。インフレーションが終わると、凍結した変動は隣接する領域をわずかに異なる時刻に再加熱させ、微視的な量子ノイズを原始密度パターンへと変換する。}

\section{どの定数が普遍的か}

我々は素粒子について極めて正確な量子論を有している。そのパラメータにはクォークとレプトンの質量、結合定数、混合データ、電荷、そして粒子そのものの一覧が含まれる。無限の計算能力があれば、この理論は加速器実験だけでなく原子、分子、化学をも説明するであろう。原理的にはその帰結は生物学や思考の基礎となる物理過程にまで及ぶ。しかし、それによってそのような系が計算しやすくなるわけではない。これは標準模型が通常の物質の微視的基礎として意図されていることを述べているにすぎない。重力は別に、古典的レベルで一般相対性理論によって記述される。両理論は強力であるが、どちらも低エネルギー物理学に現れるすべての数を必然に見せるわけではない。

陽子と電子の質量比は完全に立派な物理的数であるが、それは二つの特定の粒子に結びついている。陽子を中性子に、あるいは電子をミューオンに置き換えれば異なる比が得られる。そのようなパラメータは我々がまだ知らない深い物理学を符号化しているかもしれないが、自らの普遍性を宣伝してはいない。三つの次元定数は異なる。
\[
c,\qquad \hbar,\qquad G.
\]
それらを定義する法則には繰り返し\emph{すべての}という語が現れる。すべての物質と情報は制限速度\(c\)に従う。すべての量子系は\(\hbar\)に支配される不確定性関係に従う。すべての形態のエネルギーは同じニュートン定数\(G\)で重力を及ぼす。この広範さが、これらの定数を現在の議論において根本的たらしめている。

プランクは、量子力学の完全な意味が知られる前に、この組み合わせの重要性を認識していた。彼の定数はまず黒体放射の記述に登場した。その時点でもすでに次元的事実は明白であった。すなわち、\(c\)、\(\hbar\)、\(G\)から、いかなる特定の原子や粒子も参照することなく、長さ、時間、質量の単位を構成できるのである。

\section{プランクスケールの構成}

\(L\)、\(T\)、\(M\)がそれぞれ長さ、時間、質量の次元を表すものとする。光速度の次元は
\[
[c]=LT^{-1}.
\]
である。不確定性関係\(\Delta x\,\Delta p\gtrsim\hbar\)と\(p=mv\)を合わせると
\[
[\hbar]=ML^2T^{-1}.
\]
が得られる。最後に、ニュートンの方程式
\[
F=\frac{Gm_1m_2}{r^2},
\qquad
F=ma
\]
から
\[
[G]=L^3M^{-1}T^{-2}.
\]
が得られる。この導出が重要であるのは、次元解析が当て推量ではないからである。すなわち、各指数は不要な次元が相殺されることを要求することによって確定される。

次の形の長さを求める。
\begin{equation}
\ell=c^\alpha\hbar^\beta G^\gamma.
\label{eq:c10-length-ansatz}
\end{equation}
質量の指数は消えなければならないので、
\[
\beta-\gamma=0.
\]
時間の指数は消えなければならず、長さの指数は1に等しくなければならない。
\[
-\alpha-\beta-2\gamma=0,
\qquad
\alpha+2\beta+3\gamma=1.
\]
\(\gamma=\beta\)を用いると、これらは
\[
\alpha+3\beta=0,
\qquad
\alpha+5\beta=1.
\]
となる。その差から\(2\beta=1\)が得られ、したがって
\[
\beta=\gamma=\frac12,
\qquad
\alpha=-\frac32.
\]
となる。よって、三つの普遍定数から作られる唯一の長さは
\begin{equation}
\ell_{\mathrm P}
=\sqrt{\frac{\hbar G}{c^3}}
\approx 1.6\times10^{-35}\,\mathrm m.
\label{eq:c10-planck-length}
\end{equation}
である。講義では桁だけ、\(10^{-35}\,\mathrm m\)を保持するが、次元的導出は組み合わせそのものを確定する。

対応する時間は光が横切る時間である。
\begin{equation}
t_{\mathrm P}
=\frac{\ell_{\mathrm P}}{c}
=\sqrt{\frac{\hbar G}{c^5}}
\approx 5.4\times10^{-44}\,\mathrm s.
\label{eq:c10-planck-time}
\end{equation}
桁程度の精度ではこれは\(10^{-43}\,\mathrm s\)である。プランク長と普遍的な速度がわかれば、これは独立な構成ではない。

質量については、次元関係\(Et\sim\hbar\)と\(E=mc^2\)を用いる。\(t=t_{\mathrm P}\)とすると
\begin{equation}
m_{\mathrm P}
=\frac{\hbar}{c^2t_{\mathrm P}}
=\sqrt{\frac{\hbar c}{G}}
\approx 2.2\times10^{-8}\,\mathrm{kg}.
\label{eq:c10-planck-mass}
\end{equation}
が得られる。これは素粒子の基準からすれば驚くほど巨視的であり、およそ細かい塵の粒の質量である。その静止エネルギーは\(10^9\,\mathrm J\)程度であり、燃料タンク一杯の化学エネルギーと同程度の規模である。自動車が数回の給油で数百マイルも走れるのは、燃焼が燃料の静止エネルギーのほんのわずかな割合しか放出しないからである。一つのプランク質量を完全にエネルギーへ変換することは、まったく別の話であろう。\(c^2\)を掛ければ、微小な質量が莫大なエネルギーを表しうるのである。

加速器物理学との対比は厳しい。講義当時の高エネルギー衝突で探られる最短スケールに特徴的な\(10^{-18}\,\mathrm m\)程度の距離は、依然として約\(10^{17}\)プランク長である。したがって、電弱スケールとプランクスケールの間の未知の領域は広大である。原子核物理学が適切な類推を与える。原子核は遠くから見れば素粒子的に見えるが、より短い距離では陽子と中性子に分解され、それらはさらにクォークとグルーオンに分解される。低エネルギーの粒子は原子核に比べて単純であり得るが、それでもプランクスケールにある何かに比べれば複雑な創発的対象である。標準模型の多くのパラメータは、その階層を俯瞰すればそれほど驚くべきものではないが、観測された値と微調整は依然として説明されないままである。

\begin{classroomqa}
\audiencequestion{ヒッグス粒子は予想より軽いのか重いのか、そしてなぜその質量は不可解なのか。}
\lecturerresponse{二つの比較がある。プランク質量に対して、\(125\,\mathrm{GeV}\)のヒッグス粒子は極めて軽い。保護構造を持たない理論では、量子補正はスカラー質量をはるかに高いスケールに敏感にしがちであるため、観測値を得ることは狭い縁でパラメータを釣り合わせるように見える。当時議論されたいくつかの単純な超対称性の予想に対しては、観測されたヒッグス粒子はむしろやや重いと呼べるかもしれない。見かけの矛盾は、それを二つの異なる基準スケールと比較することから生じる。根底にある階層問題は残ったままである。\lecturetimestamp{00:26:45--00:30:54}}
\end{classroomqa}

\begin{classroomqa}
\audiencequestion{ヒッグス粒子が粒子に質量を与えるのを助けるなら、暗黒物質粒子はプランクスケール近くの質量を持つべきか。}
\lecturerresponse{いいえ。WIMPは\emph{弱く相互作用する大質量粒子}であり、この語句は実験室での検出を可能にし得る相互作用スケールを指しており、プランク質量を指しているのではない。通常のWIMPの仮定の下では、質量は陽子質量よりはるかに大きくても、\(m_{\mathrm P}\)よりは途方もなく小さいことがあり得る。通常の物質との相互作用が任意に弱くなり得るなら、宇宙論ははるかに広い質量範囲を許容するが、それに対応して直接発見は難しくなる。陽子と中性子も有用な警告を与える。それらの質量の大部分は強い相互作用による結合エネルギーであり、単にヒッグス粒子が生成したクォーク質量の和ではない。\lecturetimestamp{00:30:55--00:33:11}}
\end{classroomqa}

\begin{classroomqa}
\audiencequestion{プランク質量は何を意味するのか、そしてビット当たりの基本的なエネルギーは存在するのか。}
\lecturerresponse{有用な重力学的解釈は、ブラックホールは従来の半古典的ブラックホールであり続けながらプランクスケールよりパラメータ的に小さくすることはできないというものである。そのような対象は\(m_{\mathrm P}\)程度の質量、1プランク面積程度の事象の地平面積、1単位程度のエントロピーを持つ。その限られたブラックホールの意味では、1程度の情報にプランク程度のエネルギーを関連付けることができる。これは、あらゆる装置のあらゆる記憶ビットがプランクエネルギーを要するという普遍的な主張ではない。通常の粒子や通常の記憶ははるかに軽くあり得る。\lecturetimestamp{00:33:11--00:35:10}}
\end{classroomqa}

ここからはプランク単位を使うのが便利である。
\begin{equation}
c=\hbar=G=1.
\label{eq:c10-planck-units}
\end{equation}
長さ、時間、質量、エネルギーは、方程式~\eqref{eq:c10-planck-length}--\eqref{eq:c10-planck-mass}に対して測られた純粋な数として表すことができる。

\section{ほぼ一様なマイクロ波の空}

宇宙論の問題は全天のマイクロ波地図から始まる。WMAP、すなわちウィルキンソン・マイクロ波異方性探査機は、原始プラズマが透明になった時代の近くで放出された放射を測定した。光子はそれ以来およそ千倍の赤方偏移を受け、現在では約\(2.7\,\mathrm K\)のほぼ完全な黒体を形成している。我々が外を見るとき、共通の現在時刻における宇宙を見ているのではない。より大きな距離はより早い放出を意味し、マイクロ波背景放射は光子が初めて自由に我々に向かって流れ出した遠方の殻である。

\begin{classroomqa}
\audiencequestion{楕円形のWMAP画像は正確には何を表しているのか。}
\lecturerresponse{楕円形の各点は空のある方向を表しており、球は地図投影によって平坦化されている。左右の端は繋がり、上下の部分は反対の銀河半球を表し、銀河面は中央を横切っている。我々自身の銀河からの前景マイクロ波放射はモデル化して差し引かなければならない。残るものは近くの恒星の地図ではなく、最終散乱面におけるマイクロ波温度の方向地図である。\lecturetimestamp{00:37:49--00:40:49}}
\end{classroomqa}

第一近似では、結果は等方である：
\[
T(\hat{\mathbf n})\simeq 2.7\,\mathrm K
\]
どの視線方向\(\hat{\mathbf n}\)に対しても成り立つ。すべての点の周りで完全な等方性があれば、完全な一様性が導かれる。しかし、完全に一様な宇宙は、前回の講義で議論した銀河、銀河団、そしてコースティックのような宇宙網を自発的に生み出すことはない。何らかの非一様性が存在したに違いないが、その振幅はどんなに小さくてもよい。

方向ごとの温度のずれを次のように書く：
\[
\delta T(\hat{\mathbf n})=T(\hat{\mathbf n})-\bar T.
\]
測定された相対変動は程度である：
\begin{equation}
\frac{\delta T}{T}\sim10^{-5}.
\label{eq:c10-cmb-anisotropy}
\end{equation}
この数値への一つの道筋は力学的であった：今日の物質分布を過去に遡って進化させ、デカップリング時にどれだけの初期コントラストが必要かを問う。もう一つは直接的であった：ほとんど特徴のないマイクロ波マップを5桁増幅し、そのまだら模様自体を測定する。二つの道筋は、この議論に必要な水準で一致する。

まだら模様は多くの角度スケールで生じる。大きなパッチは小さなパッチを含み、ある領域を拡大すると統計的に類似したテクスチャが得られる。\emph{スケール不変}という用語は、すべてのスケールが文字通り同一であることを意味しない；それは、揺らぎのパワーが関連する範囲でスケールとともに弱くしか変化しないことを意味する。このほぼスケール不変性は、単なる全体の振幅ではなく、理論が説明しなければならない証拠の一部である。

なぜ温度が物質を追跡すべきなのか？重力ポテンシャルの井戸から脱出する光子はエネルギーを失う。したがって、より密度の高い領域から到来する放射は、より密度の低い領域からの放射に対してシフトされうる。講義の桁程度のレベルでは、これがデカップリング付近での
\begin{equation}
\frac{\delta\rho}{\rho}\sim10^{-5}
\label{eq:c10-density-target}
\end{equation}
を動機づける。精密なCMB計算には音響力学、重力ポテンシャル、伝達関数が含まれるので、式~\eqref{eq:c10-density-target}は\(\delta T/T\)と\(\delta\rho/\rho\)の間の厳密な等式ではなく、スケールを設定する略記である。

\begin{classroomqa}
\audiencequestion{まだら模様は最終散乱で生成された通常の微視的ノイズでありうるか？}
\lecturerresponse{観測された角度スケールではそうではない。パッチはデカップリング時点で巨大な物理的領域に対応しており、再結合プラズマ中の原子スケールの擾乱ではなかった。しかし、それらをはるかに過去に遡ると、その物理的波長は小さくなる。インフレーションが欠けていた拡大を供給する：量子的で微視的な擾乱として始まった揺らぎは、天文学的なスケールに引き伸ばされうる。\lecturetimestamp{00:46:42--00:48:02}}
\end{classroomqa}

この機構は三つの物理から理解できる：調和振動子の零点運動、振動子の減衰、そして膨張宇宙におけるスカラー場の波動方程式である。この順に組み立てよう。

\section{振動子の零点運動}

単位質量の調和振動子に対して、
\begin{equation}
E=\frac12\dot x^{\,2}+\frac12\omega^2x^2,
\qquad
\ddot x+\omega^2x=0.
\label{eq:c10-undamped}
\end{equation}
運動エネルギーとポテンシャルエネルギーは絶えず交換される。一振動にわたって平均すると、それぞれが全エネルギーの半分を担う。

量子力学はエネルギー準位を与える：
\begin{equation}
E_n=\left(n+\frac12\right)\hbar\omega.
\label{eq:c10-oscillator-levels}
\end{equation}
基底状態は\(x=0\)と\(\dot x=0\)の両方を持つことはできない。それは位置と運動量を同時に正確に指定することになるからである。その不可約なエネルギー、
\[
E_0=\frac12\hbar\omega,
\]
が零点エネルギーである。\emph{振動}という言葉は統計的に理解されるべきである：定常な基底状態は古典的な正弦軌道を描かないが、位置と運動量に非ゼロの分散を持つ。

基底状態エネルギーの半分は平均してポテンシャル項にある：
\[
\frac12\omega^2\langle x^2\rangle
=\frac14\hbar\omega.
\]
したがって
\begin{equation}
\langle x^2\rangle=\frac{\hbar}{2\omega}
\sim\frac{\hbar}{\omega}.
\label{eq:c10-zero-point-position}
\end{equation}
因子2を無視すると、揺らぎのスケールに対する複素表記は
\[
x(t)\sim\sqrt{\frac{\hbar}{\omega}}\,e^{i\omega t}.
\]
微分すると次が得られる：
\begin{equation}
\langle\dot x^{\,2}\rangle\sim\hbar\omega,
\label{eq:c10-zero-point-velocity}
\end{equation}
これは平均運動エネルギーとポテンシャルエネルギーが等しいことと整合する。冷却は\(n>0\)の励起を除去できるが、これらの基底状態の分散を除去することはできない。

\section{減衰、臨界性、凍結}

振動子を粘性媒質に浸す。蜂蜜の中を動くばねと質量は次に従う：
\begin{equation}
\ddot x+\gamma\dot x+\omega^2x=0,
\label{eq:c10-damped}
\end{equation}
ここで\(\gamma>0\)は抗力の大きさを測る。\(x=e^{\alpha t}\)を探すと次が得られる：
\begin{equation}
\alpha^2+\gamma\alpha+\omega^2=0,
\end{equation}
したがって
\begin{equation}
\alpha_\pm
=\frac{-\gamma\pm\sqrt{\gamma^2-4\omega^2}}{2}.
\label{eq:c10-damped-roots}
\end{equation}
二つの根は、二階微分方程式が要求する二つの独立な解を与える。

\begin{classroomqa}
\audiencequestion{式~\eqref{eq:c10-damped}の抗力項は負符号を持つべきか？}
\lecturerresponse{抗力自体は\(F_{\rm drag}=-\gamma\dot x\)である。ニュートンの方程式は\(\ddot x=-\gamma\dot x-\omega^2x\)である。両方の力の項を左辺に移すと、式~\eqref{eq:c10-damped}がプラス符号で得られる。したがって、微分方程式の符号と物理的な力の向きは整合している。\lecturetimestamp{00:59:38--01:00:10}}
\end{classroomqa}

\(\gamma^2<4\omega^2\)ならば、根は複素数である。次のように定義する：
\[
\Omega=\sqrt{\omega^2-\frac{\gamma^2}{4}}.
\]
すると
\begin{equation}
x(t)=e^{-\gamma t/2}
\left(A\cos\Omega t+B\sin\Omega t\right).
\label{eq:c10-underdamped}
\end{equation}
振動子は依然として平衡点を繰り返し横切るが、その包絡線は減衰する。これが不足減衰領域である。

\(\gamma^2>4\omega^2\) ならば、式~\eqref{eq:c10-damped-roots}の両根は実で負である：
\begin{equation}
x(t)=A e^{\alpha_+t}+B e^{\alpha_-t}.
\label{eq:c10-overdamped}
\end{equation}
持続的な振動はなく、両成分は減衰する。これが過減衰領域である。
\begin{equation}
\gamma=2\omega
\label{eq:c10-critical}
\end{equation}
では根が一致し、臨界減衰、すなわち二つの振る舞いの間の移行を画す。

一つの極限の場合が特に重要である。復元力がなければ、\(\omega=0\)、
\[
\alpha(\alpha+\gamma)=0,
\]
したがって
\begin{equation}
x(t)=c_0+c_1e^{-\gamma t}.
\label{eq:c10-no-restoring}
\end{equation}
速度は消えるが、座標は必ずしもゼロに戻る必要はない。それは初期条件によって選ばれる一定値 \(c_0\) に凍結する。ばねを付けずに蜂蜜の中を押し通された石は、抵抗がその運動を使い果たしたところで静止する。

ここで振動数をゆっくり変化させよう。\(\omega\gg\gamma\) から始めると、運動は振動する。減衰の尺度をほぼ固定したまま \(\omega(t)\) を減少させる。振動は遅くなり、\(\gamma=2\omega\) を横切り、過減衰となり、\(\omega\to0\) として力のない形に近づく。等価に、復元力に対して減衰を高めてもよい。宇宙論的応用では、決定的な変化は有効振動数の低下であり、ハッブル減衰はほぼ一定に留まる。最終的な凍結値は移行だけから推論することはできず、それ以前の完全な発展と初期量子状態から従う。

\section{膨張宇宙におけるスカラー波}

インフラトン \(\phi\) に戻ろう。そのポテンシャルは、高く広いプラトーと、それに続くより低い極小への急な下降を含むと仮定される。プラトー上では重要な近似は
\begin{equation}
V'(\phi)\simeq0,
\label{eq:c10-flat-potential}
\end{equation}
であり、\(V(\phi)\simeq0\) ではない。傾きは小さく、一方でエネルギー密度はインフレーションを駆動するのに十分大きいままでいられる。この形の詳細な起源は与えられない。それは模型の現象論的要請である。

\begin{figure}[htbp]
\centering
\includegraphics[width=0.88\textwidth]{lecture_10_figure_02.png}
\caption{インフラトンのプラトーと崖は左側に見えたままで、平坦空間の波動方程式が右側で始められる。黒板は近似条件 \(V'(\phi)=0\) と二階微分項の始まりを確実に支えているが、講師は未完成の方程式を隠しており、式~\eqref{eq:c10-flat-wave}は話された導出から再構成される。}
\label{fig:c10-potential-wave}
\lectureframe{01:17:17.865}
\end{figure}

平坦時空におけるほぼ自由なスカラー場に対して、
\begin{equation}
\frac{\partial^2\phi}{\partial t^2}
-\nabla^2\phi=0.
\label{eq:c10-flat-wave}
\end{equation}
非平坦なポテンシャルは右辺に \(-V'(\phi)\) を加えるであろう。プラトー近似ではそれを省く。場が完全に一様であれば、空間微分は消える。膨張宇宙では、残る一様な方程式は
\[
\ddot\phi+3H\dot\phi=0,
\]
となり、前回の講義のハッブル摩擦方程式である。

\begin{figure}[htbp]
\centering
\includegraphics[width=0.88\textwidth]{lecture_10_figure_03.png}
\caption{プラトーは左側に示され、波動方程式の successive な段階が右側に現れる。上の平坦空間演算子は読める。下の行は \(+3H\) に達するが未完成であり、したがって完全な減衰項は同時の説明から補われる。}
\label{fig:c10-hubble-board}
\lectureframe{01:18:57.554}
\end{figure}

\begin{figure}[htbp]
\centering
\begin{coursefitfigure}
\begin{tikzpicture}[x=1cm,y=0.9cm,>=Latex]
  \draw[->] (-0.2,0) -- (9.3,0) node[right] {\(\phi\)};
  \draw[->] (0,-0.2) -- (0,4.6) node[above] {\(V(\phi)\)};
  \draw[very thick,blue!55!black,smooth]
    plot coordinates {(0.25,3.9) (1.4,3.88) (2.8,3.8) (4.2,3.6) (5.3,3.25) (6.0,2.6) (6.5,1.45) (6.9,0.45) (7.5,0.28) (8.5,0.7)};
  \draw[dashed] (6.45,0) -- (6.45,2.0);
  \draw[->,thick,red!65!black] (2.0,4.25) -- (3.1,3.82);
  \node[anchor=south west] at (0.8,4.0) {slow-roll plateau: \(V'\simeq0\)};
  \node[anchor=west] at (6.55,2.05) {end of inflation};
  \node[anchor=north] at (7.45,0.22) {reheating region};
\end{tikzpicture}
\end{coursefitfigure}
\caption{議論全体で用いられるポテンシャルの明快な再構成。浅い傾きはスローロールを可能にし、\(V\) の高い値が \(H\) をほぼ一定に保つ。急な下降はインフレーションを終わらせ、蓄えられたエネルギーを解放する。}
\label{fig:c10-potential-reconstruction}
\end{figure}

非一様な場に対しては、座標距離と物理的距離を区別しなければならない。無次元の共動ラベル \(\mathbf x\) を選ぶと、無限小の固有距離は
\[
d\boldsymbol\ell=a(t)\,d\mathbf x.
\]
である。時刻を固定すると、
\[
\nabla_{\rm proper}=\frac{1}{a}\nabla_{\mathbf x},
\qquad
\nabla_{\rm proper}^2=\frac{1}{a^2}\nabla_{\mathbf x}^2.
\]
したがってスカラー方程式は
\begin{equation}
\ddot\phi+3H\dot\phi-\frac{1}{a^2}\nabla^2\phi=0,
\label{eq:c10-frw-wave}
\end{equation}
であり、ここで \(\nabla\) は共動ラベルに関して微分する。代わりに無次元のスケール因子を選ぶと、共動座標は長さを持つ。物理は変わらない。

\begin{classroomqa}
\audiencequestion{式~\eqref{eq:c10-frw-wave}の空間項は正であるべきか。}
\lecturerresponse{いいえ。平坦空間の波動演算子はすでに相対符号 \(\partial_t^2-\nabla^2\) を持つ。膨張は空間微分を因子 \(a^{-2}\) だけ変えるが、その符号を逆にはしない。フーリエ代入の後、二つの負符号が結合し、モード方程式は期待される正の復元項 \(+k^2\phi_k/a^2\) を得る。\lecturetimestamp{01:22:18--01:23:06}}
\end{classroomqa}

\section{各フーリエモードは減衰振動子である}

一度に一つの空間波を調べる：
\begin{equation}
\phi(\mathbf x,t)
=\phi_{\mathbf k}(t)e^{i\mathbf k\cdot\mathbf x}.
\label{eq:c10-mode-ansatz}
\end{equation}
複素指数はサインとコサインのための記帳である。任意の実プロファイルは、\(\mathbf k\) と \(-\mathbf k\) の間の適切な実数関係を持つモードを重ね合わせて作られる。

時間微分は \(\phi_{\mathbf k}(t)\) にのみ作用し、一方
\[
\nabla^2e^{i\mathbf k\cdot\mathbf x}
=-k^2e^{i\mathbf k\cdot\mathbf x}.
\]
式~\eqref{eq:c10-frw-wave}への代入は
\begin{equation}
\ddot\phi_{\mathbf k}
+3H\dot\phi_{\mathbf k}
+\frac{k^2}{a^2}\phi_{\mathbf k}=0.
\label{eq:c10-mode-equation}
\end{equation}
を与える。したがって各波数は
\begin{equation}
\gamma=3H,
\qquad
\omega_k(t)=\frac{k}{a(t)}.
\label{eq:c10-mode-identification}
\end{equation}
を持つ減衰振動子のように振る舞う。物理的波数 \(k/a\) は宇宙の膨張に伴って赤方偏移する。その結果、有効復元力 \(k^2/a^2\) はゼロに向かって減少する。

短い物理的波長で始まるモードは\(\omega_k\gg H\)を持ち、振動する。その波長は引き伸ばされ、振動数は低下し、やがてハッブル減衰が支配的になる。するとそれはゼロに引き戻されるのではなく、振動しない値に近づく。異なる\(\mathbf k\)モードは異なる時刻にクロスオーバーする。より大きな\(k\)は、固定した\(a\)のもとではより短い波長に対応し、その物理的振動数がハッブルスケールまで低下するには、より多くの膨張を待たなければならない。

\section{地平線退出と凍結した真空揺らぎ}

振動子の臨界条件をオーダー1の見積もりとして用いると、
\[
\gamma=2\omega_k,
\]
式~\eqref{eq:c10-mode-identification}より
\begin{equation}
3H=\frac{2k}{a},
\qquad
a_{\rm cross}=\frac{2k}{3H}.
\label{eq:c10-critical-mode}
\end{equation}
が得られる。数値因子は概念的要点ではない。クロスオーバーは、物理的波数\(k/a\)が\(H\)のオーダーになるときに起こる。

次を定義する。
\[
\lambda_{\rm com}=\frac{2\pi}{k},
\qquad
\lambda_{\rm phys}=a\lambda_{\rm com}
=\frac{2\pi a}{k}.
\]
次元解析により\(a\)の配置が決まる。無次元の\(\mathbf x\)のもとで、\(k\)と\(\mathbf k\cdot\mathbf x\)はともに無次元であり、一方\(a\)と\(\lambda_{\rm phys}\)は長さを持つ。したがって式~\eqref{eq:c10-critical-mode}は
\begin{equation}
\lambda_{\rm phys}\sim H^{-1},
\label{eq:c10-horizon-crossing}
\end{equation}
を、\(2\)、\(3\)、\(2\pi\)のような因子を除いて意味する。

ハッブル則\(v=Hd\)は距離\(H^{-1}\)を\(c=1\)のオーダーの後退速度と結びつける。インフレーションの間、\(H\)は大きく、ほぼ一定であるため、このハッブル長は因果的スケールの役割を果たすにもかかわらず微視的である。この出来事を\emph{地平線退出}と呼ぶのは慣習的である。物質的な物体が物理的な壁を通り抜けるわけではない。モードの物理的波長はほぼ固定されたハッブルスケールを超えて成長し、その力学は急速な振動からほとんど凍結した振幅へと変化する。

\begin{figure}[htbp]
\centering
\begin{coursefitfigure}
\begin{tikzpicture}[x=1cm,y=0.85cm,>=Latex]
  \draw[->] (0,0) -- (9.2,0) node[right] {cosmic time};
  \draw[->] (0,0) -- (0,5.0) node[above] {physical length};
  \draw[dashed,thick] (0.25,2.75) -- (8.75,2.75)
    node[right] {\(H^{-1}\)};
  \draw[very thick,blue!55!black]
    (0.35,0.55) .. controls (3.3,0.75) and (5.3,1.55) .. (8.5,4.55);
  \node[blue!55!black,anchor=west] at (6.65,4.55) {\(\lambda_{\rm phys}=2\pi a/k\)};
  \draw[densely dotted] (6.05,0) -- (6.05,2.75);
  \fill (6.05,2.75) circle (2pt);
  \node[anchor=north] at (6.05,0) {horizon exit};
  \node[align=center] at (2.6,1.7) {oscillating\\sub-Hubble mode};
\end{tikzpicture}
\end{coursefitfigure}
\caption{ほぼ指数関数的な膨張の間、物理的波長は成長する一方でハッブル長はほぼ固定されたままである。両者の交差は、オーダー1の因子を除いて、急速な振動から凍結したモードへの移行を印づける。}
\label{fig:c10-horizon-reconstruction}
\end{figure}

真空揺らぎのため、インフレーションがすべての波を単に消し去ることはない。すべての短波長モードは零点運動から始まる。膨張は一つのモードをハッブルスケールまで引き伸ばして凍結させるが、さらに短い真空モードは振動を続ける。それらも順に引き伸ばされて凍結する。この繰り返しにより、広い範囲の波長と向きにわたる凍結した寄与を含む場のプロファイルが構築される。結果は確率的であるが、構造を持たないわけではない。その統計は量子状態とインフレーション背景から従う。

\begin{figure}[htbp]
\centering
\includegraphics[width=0.88\textwidth]{lecture_10_figure_04.png}
\caption{長い場の波と短い場の波の盤面全体にわたる重ね合わせであり、ラベル\(\gamma\)と\(3H\)のそばに小さな減衰のスケッチがある。この図は定性的なモード像を記録するものであり、厳密な関数を読み取るために用いられるものではない。}
\label{fig:c10-mode-board}
\lectureframe{01:45:28.476}
\end{figure}

\begin{classroomqa}
\audiencequestion{モードは一点で凍結するのか、それとも波長で凍結するのか、そして\(H\)自体が指数関数的に成長しているのか。}
\lecturerresponse{それらは波長がハッブルスケールに達したときに凍結する。ほぼ平坦なプラトー上では、\(V\)、したがって\(H\)はほぼ一定である。指数関数的に成長するのはスケール因子\(a(t)\)であり、それゆえ物理的波長\(a/k\)である。各共動モードはそれぞれの時刻に条件\(k/a\sim H\)に達する。\lecturetimestamp{01:45:06--01:45:39}}
\end{classroomqa}

\begin{classroomqa}
\audiencequestion{これらは電磁波なのか。}
\lecturerresponse{いいえ。これらはインフラトン・スカラー場の揺らぎである。電磁真空揺らぎはよく知られた類推を与えるが、式~\eqref{eq:c10-mode-equation}の変数は\(\phi\)である。\lecturetimestamp{01:45:39--01:46:04}}
\end{classroomqa}

\begin{classroomqa}
\audiencequestion{波長の連続体はあるのか。}
\lecturerresponse{無限の空間という理想化に対しては、はい。場はフーリエ積分
\begin{equation}
\phi(\mathbf x,t)
=\int\frac{d^3k}{(2\pi)^3}\,
\phi_{\mathbf k}(t)e^{i\mathbf k\cdot\mathbf x}.
\label{eq:c10-fourier-integral}
\end{equation}
で表される。すべての波数が同じ赤方偏移と凍結の機構を経る。\lecturetimestamp{01:46:04--01:46:34}}
\end{classroomqa}

\begin{classroomqa}
\audiencequestion{有限の宇宙では離散的なモードしかないのか。}
\lecturerresponse{はい。有限の空間体積上の境界条件は許される波数を量子化する。議論しているスケールよりずっと大きな体積に対しては、間隔は非常に細かくなり得て、和は式~\eqref{eq:c10-fourier-integral}の積分でよく近似される。連続的な描像は理想化であり、異なる局所的機構ではない。\lecturetimestamp{01:46:34--01:47:01}}
\end{classroomqa}

ここで二つの運動が共存する。空間平均された場\(\bar\phi(t)\)はポテンシャルに沿ってゆっくりと古典的に転がる。その上に小さな量子揺らぎ
\[
\phi(\mathbf x,t)=\bar\phi(t)+\delta\phi(\mathbf x,t),
\]
が、凍結したモードから組み立てられて乗っている。揺らぎのエネルギー自体は、観測される物質のコントラストとなるには小さすぎる。その重要性は、インフレーションが終わる時刻を摂動することにある。

完全な単一場の量子計算は、この振動子による議論をより精密にする。軽いスカラー場は、波数の対数区間あたり、程度の典型的な揺らぎを獲得する。
\begin{equation}
\delta\phi_{\rm exit}\sim\frac{H}{2\pi}.
\label{eq:c10-field-amplitude}
\end{equation}
規格化の導出には、インフレーション背景における場のモードの量子化が必要である。振動子の描像はすでに本質的な機構を与えている。すなわち、真空の揺らぎは基底状態においてさえ存在し、膨張がそれらを古典的に見える凍結したパターンへと変換するのである。

\section{崖は場のノイズを密度のノイズへ変換する}

まず\(\delta\phi\)を抑える。任意の固定した時刻において、場はどこでも一つの値をもち、空間\(x\)を上向きに、場の値\(\phi\)を横向きに描くとき、黒板では垂直な線で表される。時間が進むにつれて、その線は崖へ向かって動く。宇宙全体が一緒に落下点へ到達し、インフラトンのポテンシャルエネルギーは粒子、放射、熱となり、指数関数的膨張は終わり、その結果生じたエネルギー密度は通常の膨張とともに薄まり始める。

\begin{figure}[htbp]
\centering
\includegraphics[width=0.88\textwidth]{lecture_10_figure_05.png}
\caption{黒板の空間対場の図。縦軸は\(x\)、横方向は\(\phi\)であり、かすかなほぼ垂直のプロファイルが successive な場の配置を示し、右側の斜線を施した境界がインフレーションの終わる崖である。}
\label{fig:c10-cliff-board}
\lectureframe{01:50:54.502}
\end{figure}

\begin{figure}[htbp]
\centering
\begin{coursefitfigure}
\begin{tikzpicture}[x=1cm,y=0.82cm,>=Latex]
  \draw[->] (0,0) -- (9.2,0) node[right] {\(\phi\)};
  \draw[->] (0,0) -- (0,5.2) node[above] {space \(x\)};
  \draw[thick,gray!60] (8.2,0.15) rectangle (8.55,4.85);
  \foreach \y in {0.35,0.85,...,4.55} {
    \draw[gray!60] (8.2,\y) -- (8.55,{\y+0.25});
  }
  \node[anchor=south,rotate=90] at (8.38,2.5) {cliff};
  \draw[thick,densely dashed,gray!70] (3.2,0.2) -- (3.2,4.75);
  \draw[thick,blue!45!black,smooth]
    plot coordinates {(5.0,0.2) (5.12,0.75) (4.9,1.25) (5.18,1.8) (4.94,2.35) (5.2,2.9) (4.88,3.45) (5.12,4.05) (5.0,4.75)};
  \draw[very thick,red!65!black,smooth]
    plot coordinates {(7.15,0.2) (7.32,0.7) (7.0,1.2) (7.42,1.75) (7.1,2.3) (7.5,2.85) (7.02,3.35) (7.35,4.0) (7.16,4.75)};
  \draw[->,thick] (3.55,2.5) -- (4.55,2.5);
  \draw[->,thick] (5.55,2.5) -- (6.65,2.5);
  \node[align=center,anchor=south,font=\small] at (3.2,4.8) {uniform\\profile};
  \node[align=center,anchor=south,font=\small] at (5.0,4.8) {frozen\\ripple};
  \node[align=center,anchor=south,font=\small] at (7.2,4.8) {near\\reheating};
\end{tikzpicture}
\end{coursefitfigure}
\caption{空間対場の描像の明快な再構成。平均場は崖へ向かって動く一方で、凍結した空間的なさざ波が持続するため、異なる位置はわずかに異なる時刻にインフレーションの終わりに到達する。}
\label{fig:c10-cliff-reconstruction}
\end{figure}

凍結した変動を復元する。線は今やわずかに波打っている。隣り合う場所は\(\phi\)のわずかに異なる値をもつ。肩を並べて崖へ行進しようとする人々の列を考えよ。数人は少し先に進み、数人は少し後れている。先に進んだ者たちが最初に崖を越える。

このタイミングのずれこそが密度コントラストの真の源である。先に越える領域は先に再加熱する。そのポテンシャルエネルギーは粒子と熱に注ぎ込まれ、その後、そのエネルギーは宇宙がインフレーション後の時代に膨張する間に薄まる時間をもつ。わずかに遅れて越える近傍の領域は遅れて再加熱するので、両者を比較すると、遅い領域の方が単に薄まる時間が短かったという理由だけでより高密度になる。したがって、最初の一連の越境は比較的高密度のポケットを残す。それらのポケットが膨張する間に、より短い波長のさざ波が別の一連の位置を縁の向こうへ運び、より大きなパターンの中により小さなポケットを作る。これを凍結したモードの階層全体にわたって繰り返すことで、構造の中に構造が生じる。したがって、微小な場の揺らぎはそれ自身のエネルギーによって重要なのではなく、再加熱が起こる局所的な時刻をずらすことによって重要なのだ。

タイミングの関係は模式的に書くことができる。一様な場が終端値\(\phi_{\rm end}\)に速度\(\dot{\bar\phi}\)で到達するなら、小さな場の変位は局所的な終了時刻を次だけ変える。
\begin{equation}
\delta t_{\rm end}
\simeq-\frac{\delta\phi}{\dot{\bar\phi}}.
\label{eq:c10-time-delay}
\end{equation}
より遅い転がりは、同じ\(\delta\phi\)をより大きな遅延に対応させる。その余分な時間の間に、より早い領域はより多く膨張し薄まる。1程度の転送因子を除いて、曲率すなわち密度の摂動はしたがって次によって支配される。
\begin{equation}
\mathcal R
\sim H\,\delta t_{\rm end}
\sim-\frac{H\,\delta\phi}{\dot{\bar\phi}}.
\label{eq:c10-curvature-estimate}
\end{equation}
これを式~\eqref{eq:c10-field-amplitude}と組み合わせると、おなじみのスローロールのスケールが得られる。
\begin{equation}
\mathcal P_{\mathcal R}^{1/2}
\sim\frac{H^2}{2\pi|\dot{\bar\phi}|},
\label{eq:c10-slow-roll-amplitude}
\end{equation}
正確な規約と、その後の\(\delta\rho/\rho\)への転送には、より完全な摂動計算が必要である。場の転がりが遅いほど、隣り合う崖の越境の間の時間は長くなり、その間の薄まりは大きくなる。したがって、同じ場の揺らぎは、\(\lvert\dot{\bar\phi}\rvert\)が小さいほど、より大きな密度の刻印を残す。

多くのモードが異なる時刻にハッブルスケールを横切るので、多くの異なる大きさの領域がわずかに異なる時刻に再加熱する。その結果生じる密度場は、スケールをまたいで入れ子になった高密度と低密度のポケットを含む。その統計はインフレーションのポテンシャルの高さ、傾き、曲率に依存する。適切なスローロール模型は、次に近い、ほぼスケール不変な振幅を自然に生み出す。
\[
\frac{\delta\rho}{\rho}\sim10^{-5},
\]
そして詳細なCMBスペクトルは、少数の宇宙論的パラメータを多くの角度スケールに対して検証することを可能にする。

\section{原始の種から空へ戻る}

ここで、過去の光の殻の集合の中心にいる観測者に話を戻そう。最終散乱球は三次元の原始密度場を横切っている。その殻が異なる重力ポテンシャルと交差する方向では、わずかに異なるマイクロ波温度が生じる。WMAPはその球面上の二次元パターンを表示する。この地図は時刻ゼロの写真でも、球の内部にあるすべての構造の投影でもない。マイクロ波光子が最後に散乱された特定の殻上の温度構造なのである。

殻の内部では、同じ原始揺らぎが進化し続ける。重力が過密度を増幅し、コースティックのような流れがシートやフィラメントを鋭くし、ガスが集まり、銀河が形成される。したがって、CMBのパターンとその後の宇宙網は一つの歴史の二つの段階である。凍結したインフレーション揺らぎがまず微小な密度の種として現れ、後に可視的な構造となるのである。

\begin{classroomqa}
\audiencequestion{今日私たちが目にするマイクロ波背景放射は、どの時代を表しているのか。}
\lecturerresponse{光子はおよそ\(13\)億年を旅してきたが、それらが放たれたのは熱いビッグバンの数百万年後であり、電子と陽子が中性水素に結合して宇宙が透明になったときである。インフレーションが終わったのはそれよりはるかに早い。再加熱が熱い物質と放射を生み出し、その後宇宙は膨張して冷え、最終散乱に至った。この議論にとって頑健なのは、正確な再結合の年代ではなく、桁の程度---数百万年---である。\lecturetimestamp{02:00:38--02:02:19}}
\end{classroomqa}

\section{一つの空と宇宙分散}

\begin{classroomqa}
\audiencequestion{インフレーションは、WMAPで報告された優先方向や北--南の非対称性を予言するのか。}
\lecturerresponse{ここで展開した単純なインフレーション像から、そのような優先方向は導かれない。報告された大角度の異常の位置づけは不確かであり、前景放射と統計的揺らぎは現実の懸念であり、責任ある評価にはセンセーショナルな要約ではなく専門文献が必要であった。インフレーションが予言するのは集団の統計的性質であり、一つの珍しい空の特徴が新しい物理かどうかを決めるのはより難しい。\lecturetimestamp{02:02:19--02:03:26}}
\end{classroomqa}

小さな角度スケールの特徴は空の至る所に何度も現れる。それらは多くのサンプルを与えるので、その分布と等方性を統計的に検証できる。四重極程度の大きさの特徴は異なる。観測可能な球面には、独立な大角度モードがわずかしか収まらない。したがって、基礎となる理論が正しくても、測定されるパワーは集団平均から大きくずれうる。この不可避のサンプリング不確かさが\emph{宇宙分散}である。

\begin{classroomqa}
\audiencequestion{いくつかの低い多重極の見かけの整列についてはどうか。}
\lecturerresponse{それは手がかりかもしれないし、前景放射や解析上の人為効果かもしれないし、ありそうにないが許される実現かもしれない。困難なのは、同じ統計的集団から別の観測可能な宇宙を引き出せないことである。六回のコイントスで四回表が出る硬貨が偏っているとは限らない。繰り返しの組が決めるだろうが、宇宙論が最大スケールで与えるのは一組だけである。この異常は研究に値するかもしれないが、一つの空だけでは決定的な有意性検定に必要な繰り返し試行を与えることはできない。\lecturetimestamp{02:05:31--02:07:36}}
\end{classroomqa}

\section{量子運動から宇宙構造へ}

三つの普遍定数がプランクスケールを定義し、身近な素粒子物理学が量子重力からどれほど遠いかを明らかにする。マイクロ波の空は精密な宇宙論的目標を与える。ほぼ等方であり、揺らぎの振幅はおよそ\(10^{-5}\)であり、構造は多くのスケールに広がっている。量子力学により、あらゆる場のモードは零点揺らぎをもつ。インフレーションは各インフラトンモードを、その物理的波長が\(H^{-1}\)程度になるまで引き伸ばし、その時点でハッブル減衰がその振動を抑え、実質的に凍結した振幅を残す。

一様なインフラトンは、その凍結したノイズの下で転がり続ける。場がそのプラトーの端に達すると、隣接する領域がまったく同時に再加熱するわけではない。より早い領域はより長く希薄化し、より遅い領域はより高密度のままである。これが\(\delta\phi\)を原始的な密度と曲率のパターンに変換する。その後、重力がそのパターンを銀河と宇宙網へと増幅し、最終散乱殻はその以前の状態をマイクロ波の空に保存する。したがって、微視的な零点運動、宇宙膨張、ずれた再加熱は別々の物語ではない。これらが相まって、インフレーション宇宙がなぜ並外れて滑らかでありながら、私たちが存在できるほどに構造をもつのかを説明する。
